\scanpage{9}
\begin{flushright}
\handwritten{16 July .}\\
\end{flushright}
Dear Grothendieck,\\

Here is the new version of \handwritten{4}.\handwritten{10}.\handwritten{6}.

On rereading number\textsuperscript{o} \handwritten{9} of the Seminar on Algebraic Geometry \handwritten{4} IV I found the following trifles :

page \handwritten{149} line \handwritten{16} read \(E'\handwritten{\vee} E'' = E\)

page \handwritten{158} line \handwritten{5} read \(s\varepsilon\operatorname{Hom}(e_E,G)\)

page \handwritten{158} complete the diagram \handwritten{9}.\handwritten{3}.\handwritten{6}.\handwritten{2}

page \handwritten{159} line \handwritten{15} read ``the functor \(\handwritten{\sigma}\) composed of

In numbers\textsuperscript{o} \handwritten{10} and following I found

page \handwritten{2} line -\handwritten{7} read ``hence the assertion ''

page \handwritten{14} line \handwritten{5} read \(\handwritten{\mathscr{H}}\mathrm{om}_A\)

page \handwritten{14} line \handwritten{17} ditto .

% Left margin: the two sloping underlined words stand to the left of the vertical paragraph mark.
\reversemarginpar
\margin{\handwritten{\figure{\begin{tikzpicture}[x=.85cm,y=.85cm]
\node[rotate=36,anchor=base west,inner sep=0pt] at (.15,.82) {\uline{ask}};
\node[rotate=36,anchor=base west,inner sep=0pt] at (.35,.35) {\uline{Messing}};
\draw (2.26,.35) -- (2.26,1.62);
\end{tikzpicture}}}}
\normalmarginpar
Finally , an unpleasant surprise , I noticed that number\unsure{\deleted{?}\handwritten{\textsuperscript{o}}|\handwritten{\textsuperscript{o}}} \handwritten{14} was missing from the notes you sent me (Modules on a topos defined by gluing ). I remember having written this section and I think I did give it to you . In any case I do not have it among my papers . I hope this section has not been lost .

\begin{flushright}
Yours\\
\handwritten{J. \unsure{L|V}}\\
\end{flushright}
% archive sheet number: 8
